\section{附录：术语索引与复习路线}

\subsection{术语消化：本期关键词索引}

\noindent\begin{longtable}{p{0.23\textwidth}p{0.32\textwidth}p{0.36\textwidth}}
\toprule
术语 & 一句话解释 & 在本期中的作用 \\
\midrule
OpenClaw & 触发 Agent 框架共识的产品/开源框架 & 代表 Agent 范式被产品化展示。 \\
Post-train & 预训练之后围绕任务、偏好、环境和反馈继续训练 & Agent 时代的主战场。 \\
RL Infra & 支撑 rollout、reward、policy update 和 eval 的基础设施 & 决定 Agent 后训练能否 scale。 \\
Rollout & 在任务环境中生成轨迹和交互过程 & RL 数据和反馈的来源。 \\
MiMo-V2 & 小米模型家族的访谈核心案例 & 展示多模型如何被 Agent 框架编排。 \\
1T 入场券 & 对标前沿 Agent 能力所需的基础模型规模门槛 & 说明基座模型仍重要。 \\
组织平权 & 多职能进入同一研究闭环 & 解释为什么后训练需要 diversity。 \\
环境比经验重要 & 研究速度来自算力、infra、组织和反馈环境 & 结尾的战略判断。 \\
\bottomrule
\end{longtable}

\subsection{复习路线}

建议先读第 2 章理解 Chat 到 Agent 的范式迁移，再读第 3 章掌握后训练闭环，然后读第 4--6 章理解 MiMo 编排、算力分配和组织平权。第 7--8 章适合作为产业判断：为什么 2026 年竞争不再只是基座模型，而是架构、后训练、环境和组织能力的组合。

\begin{knowledgebox}{与 EP139 的连接}
EP139 苏煜访谈给出 Agent 技术史和 universal digital agent 框架；EP138 罗福莉访谈则展示一家模型实验室如何把这个框架转化为后训练、模型编排、资源配置和组织重构。两期应连读。
\end{knowledgebox}

\subsection{本章小结}

本期关键词可以压缩为六个：OpenClaw、后训练、RL infra、多模型编排、算力重配、组织平权。它们共同指向一个判断：Agent 时代的竞争是系统竞争。


